\section{Improving constraints: combination with CMB and BAO data}
\label{sec:5.3}

\begin{figure}
\centering
    \staticfigure{figures/chapter5/fig5.10.png}{$68.3\%$, $95.4\%$ and $99.7\%$ confidence level contours on $w_1$ (here $w_a$) and $w_0$ for a flat Universe. The Union SN set was combined with CMB or BAO constraints (left) and the SN+CMB+BAO constraints with and without the SNe systematics removed (right). The diagonal line represents $w_0 + w_a = 0$: the likelihoods based on observational data remain below it, favouring matter domination at $z \gg 1$.  Image credit: Fig. 16 in \cite{Kowalski2008}.}
\end{figure}

\begin{figure}
\centering
    \staticfigure{figures/chapter5/fig5.11.png}{Comparison of $(\Omega_M , w_0)$ credible regions from the Constitution data when two different light curve fitters are used to find the luminosity of the SNe. Left: SALT; right: MCLS2k2. Notice the shift in the position of the $\Lambda$CDM parameter values (the dot) with no change in the SNe observed, merely the data calibration. Image credit: Fig. 4 from \cite{Hicken2009}.}
\end{figure}


Tightening of cosmological parameter constraints can be achieved by combining SNe data with a range of other methods for estimating the parameters, foremost amongst which are multipoles in the power spectra from the cosmological microwave background (\cmb{}) \cite{LewisBridle2002,Spergel2007,Komatsu2010,Dunkley2009} and the correlation func- tion for luminous red galaxies from baryon acoustic oscillation (\bao{}) measurements \cite{Eisenstein2005,Carroll2001,Percival2007}. These measurements are sometimes referred to as ”standard rulers” because they directly measure fundamental cosmological distances \footnote{The ”standard candle” measures the luminosity distance because the luminosity of the object is independent of its local environment, so the ratio of the (corresponding) distance between two standard candles depends only on the intervening spatial curvature (given by the FLRW metric). Similarly, the ”standard ruler” measures the angular diameter distance because the ratios of angle subtended on the sky depends only on the spatial curvature via the FLRW metric. \cite{Hobson2006}}. 

A full estimate of the cosmological parameters from such sources requires a similar process\footnote{A summary of the \bao{} process is provided in \cite{vanOsten2006}§2.; for the \cmb{} an example of parameter extraction is \cite{LewisBridle2002} while the parameters extracted from this are summarised in \cite{vanOsten2006,Carroll2001}.} to this project: generating the relevant observational data from different values of the cosmological parameters and using a Markov chain Monte Carlo algorithm to move through the parameter space based on an acceptance/rejection scheme (e.g. Metropolis-Hastings, although much effort has been put into generating efficient schemes purely for the \cmb{} analysis \cite{Christensen2001,Dunkley2005,LewisBridle2002} ). Although credible regions for the parameters $(\Omega_m , w_0 , w_1 )$ are available from the \cmb{} \cite{Spergel2007,Komatsu2010,Dunkley2009} and $(\Omega_m , w_0)$ from \bao{} measurements \cite{Eisenstein2005,Percival2007}, they are not in a form which enabled me to compare them directly to my own results (except by sight).

A trend throughout the results has been tension between the parameter estimation produced by the Gold set and that produced by the Union and Constitution sets, which is emphasised by the lack of overlap in the credible regions for the Gold 0.2x set and the Constitution/Union 0.2x sets. This lack is easily explained if one considers the formation of the projected SNAP sets. The (unknown) parameter values for our Universe would produce a $\mu(z)$ curve according to the Friedmann equations \cref{eq:1.4}, the luminosity distance formula \cref{eq:1.15} and the theoretical distance modululs formula \cref{eq:1.13}. The $\mu(z)$ points from observations of Type Ia SNe must be calculated from the pragmatic distance modulus formula \cref{eq:1.14} once their luminosities (hence luminosity distances) are found from a light-curve fitter: the choice of corrections to obtain \cref{eq:1.13} from \cref{eq:1.14} and the choice of light curve fitter to determine \cref{eq:1.15} create a particular scatter of the data points around the \lq\lq{}true\rq\rq{} $\mu(z)$ curve. Uncertainties in these choices and limitations on the instrument precision and calibration create error bars about the data points. In this way, the jiggled Snap and Snap 0.2x sets were generated from a known set of cosmological values. The 0.2x sets for Constitution, Gold and Union have reduction in the error bars, but not in the scattering (for reasons detailed in \cref{sec:5.1}). Therefore the reduced error bars shrink the subset of parameter values which provide good fits to the data, while the same data points maintain the different scatters from the different data sets. This exaggerates the difference between the data sets, so overlap becomes more unlikely as the errors shrink. The difference between the data sets results mostly from the light curve fitters used to determine the SN luminosity: the Gold set uses the MLCS2k2 fitter \cite{Riess2007}, whereas the Union set uses the SALT fitter \cite{Kowalski2008}. The Constitution set was calibrated with four different fitters (SALT being the preferred fitter): their comparison of the effects of SALT and MLCS2k2 on the parameter estimates (\cref{fig:fig5.1a-f}) shows the same trend as my results, with the MLCS2k2 fitter placing the $\Lambda$CDM parameter values within 2$\sigma$ compared to 1$\sigma$ for SALT. The shift in credible regions increases when the SNe compared are different as well as the light-curve fitters. This illustrates the vital importance of developing new light-curve fitters which cause the credible regions from different fitters to converge. Until this task is completed, there will be uncertainties in the values of the cosmological parameters derived from SNe (see §C for more detail on this). The physical values of the parameters are obscured by our lack of knowledge of both the physical mechanisms associated with Type Ia SNe and the effect of the intervening environment \cite{Leibundgut2009}.
